\documentclass[10pt,aspectratio=169]{ntuabeamer}

\graphicspath{{latex/gnns/assets/}{assets/}}

\defineslidesource{distill2021}{Sanchez-Lengeling, Reif, Pearce, Wiltschko. \emph{A Gentle Introduction to Graph Neural Networks}. Distill, 2021.}
\defineslidesource{kipf2017}{Kipf, Welling. \emph{Semi-Supervised Classification with Graph Convolutional Networks}. ICLR, 2017.}
\defineslidesource{hamilton2017}{Hamilton, Ying, Leskovec. \emph{Inductive Representation Learning on Large Graphs}. NeurIPS, 2017.}
\defineslidesource{velickovic2018}{Velickovic et al. \emph{Graph Attention Networks}. ICLR, 2018.}

\title{Graph Neural Networks}
\subtitle{A concise introduction to graph representation learning and message passing}
\author{NTUA School of Civil Engineering}
\date{\today}

\begin{document}

\begin{frame}[plain]
  \makepresentationtitle
\end{frame}

\sectiondivider{Notation and Tasks}

% TODO: Add a toy graph diagram with 5--6 nodes, node features, and a highlighted neighborhood.
\begin{frame}{Notation}
  \small
  \begin{itemize}
    \item Graph: \(G = (V, E)\), with \(N = |V|\) nodes
    \item Node features: \(\bd{X} \in \mathbb{R}^{N \times F}\), and \(\bd{H}^{(0)} = \bd{X}\)
    \item Node state vs matrix state: \(\bd{h}_{v}^{(k)}\) is one node, \(\bd{H}^{(k)}\) stacks all nodes
    \item Adjacency: \(A \in \mathbb{R}^{N \times N}\)
    \item Self-loops: \(\hat{A} = A + I\) when a layer keeps each node's own state
    \item Degree: \(\hat{d}_v = \sum_{u=1}^{N} \hat{A}_{vu}\), \(\hat{D} = \operatorname{diag}(\hat{d}_1, \dots, \hat{d}_N)\)
    \item Neighborhood: \(\hat{\mathcal N}(v) = \{u : \hat{A}_{vu} \neq 0\}\)
    \item Edge values: \(\hat{A}_{vu} \in \{0, 1\}\) or \(\hat{A}_{vu} \in \mathbb{R}\)
  \end{itemize}
\end{frame}

% TODO: Add a three-panel visual showing node-level, edge-level, and graph-level prediction examples.
\begin{frame}{Tasks}
  \textcolor{ntuablue}{Node-level}
  \begin{itemize}
    \item Citation node classification
    \item Community or role prediction
  \end{itemize}

  \textcolor{ntuablue}{Edge-level}
  \begin{itemize}
    \item Link prediction
    \item Relation prediction
  \end{itemize}

  \textcolor{ntuablue}{Graph-level}
  \begin{itemize}
    \item Molecular property prediction
    \item Whole-graph regression
  \end{itemize}
\end{frame}

\sectiondivider{Message Passing}

% TODO: Add a pipeline diagram: encoder -> L message-passing layers -> task-specific decoder.
\begin{frame}{Common GNN Architecture}
  \small
  \begin{itemize}
    \item Encode raw inputs into hidden states
    \[
      \bd{H}^{(0)} = \operatorname{ENC}_{V}(\bd{X}),
      \qquad
      \bd{E}^{(0)} = \operatorname{ENC}_{E}\!\left(\bd{E}_{\mathrm{raw}}\right)
    \]
    \item Apply \(L\) message-passing layers
    \[
      \bd{H}^{(k+1)} = f^{(k)}\!\left(\bd{H}^{(k)}, \hat{A}\right),
      \qquad k = 0, \dots, L-1
    \]
    \item Decode final embeddings for the target level
    \[
      \hat{y}_{v} = \operatorname{DEC}_{\mathrm{node}}\!\left(\bd{h}_{v}^{(L)}\right),
      \qquad
      \hat{y}_{uv} =
      \operatorname{DEC}_{\mathrm{edge}}\!\left(\bd{h}_{u}^{(L)}, \bd{h}_{v}^{(L)}\right)
    \]
    \[
      \bd{h}_{G} = \operatorname{READOUT}\!\left(\{\bd{h}_{v}^{(L)} : v \in V\}\right),
      \qquad
      \hat{y}_{G} = \operatorname{DEC}_{\mathrm{graph}}\!\left(\bd{h}_{G}\right)
    \]
  \end{itemize}
\end{frame}

\begin{frame}{Message-Passing Template}
  \textcolor{ntuablue}{Node-wise view}
  \[
    \bd{h}_{v}^{(k+1)} =
    \sigma\!\left(
      \operatorname{AGG}_{v}^{(k)}\!\left(\{\bd{h}_{u}^{(k)} : u \in \hat{\mathcal N}(v)\}\right)
      \bd{W}^{(k)}
    \right)
  \]

  \textcolor{ntuablue}{Or in matrix notation}
  \[
    \bd{H}^{(k+1)} =
    \sigma\!\left(P^{(k)} \bd{H}^{(k)} \bd{W}^{(k)}\right)
  \]

  \begin{itemize}
    \item Same message-passing layer, two views
    \item Node-wise: update one node \(v\)
    \item Matrix notation: update all nodes together
    \item \(P^{(k)}\): message-passing or aggregation operator
    \item \(\bd{W}^{(k)}\): learnable linear map
    \item GNN variants differ mainly in how they define the aggregation weights
  \end{itemize}
\end{frame}

\sectiondivider{Message-Passing Variants}

\begin{frame}{Message-Passing Variants}
  \begin{itemize}
    \item All variants update node states by aggregating neighbor information
    \item They differ in how the aggregation weights are defined
    \item Node-wise and matrix notation describe the same layer
  \end{itemize}

  \[
    \bd{H}^{(k+1)} =
    \sigma\!\left(P^{(k)} \bd{H}^{(k)} \bd{W}^{(k)}\right)
  \]
\end{frame}

% TODO: Add a small neighborhood-averaging sketch with normalized incoming edge weights.
\begin{frame}{Message-Passing Variants: Weighted Averaging}
  \begin{itemize}
    \item Fixed row-normalized weights
  \end{itemize}
  \[
    P = \hat{D}^{-1}\hat{A}
  \]

  \textcolor{ntuablue}{Node-wise update}
  \[
    \bd{h}_{v}^{(k+1)} =
    \sigma\!\left(
      \sum_{u \in \hat{\mathcal N}(v)}
      \frac{\hat{A}_{vu}}{\hat{d}_{v}} \,
      \bd{h}_{u}^{(k)} \bd{W}^{(k)}
    \right)
  \]

  \textcolor{ntuablue}{Or in matrix notation}
  \[
    \bd{H}^{(k+1)} =
    \sigma\!\left(
      \hat{D}^{-1}\hat{A}\,\bd{H}^{(k)} \bd{W}^{(k)}
    \right)
  \]

  \begin{itemize}
    \item \(\hat{A}_{vu} / \hat{d}_v\): weight assigned to neighbor \(u\)
    \item Works for binary or weighted edges
  \end{itemize}
\end{frame}

% TODO: Add a comparison visual between row-normalized propagation and symmetric normalization.
\begin{frame}{Message-Passing Variants: GCN\slideref{kipf2017}}
  \begin{itemize}
    \item Symmetric normalization
  \end{itemize}
  \[
    P = \hat{D}^{-1/2}\hat{A}\hat{D}^{-1/2}
  \]

  \textcolor{ntuablue}{Node-wise update}
  \[
    \bd{h}_{v}^{(k+1)} =
    \sigma\!\left(
      \sum_{u \in \hat{\mathcal N}(v)}
      \frac{\hat{A}_{vu}}{\sqrt{\hat{d}_{v}\hat{d}_{u}}} \,
      \bd{h}_{u}^{(k)} \bd{W}^{(k)}
    \right)
  \]

  \textcolor{ntuablue}{Or in matrix notation}
  \[
    \bd{H}^{(k+1)} =
    \sigma\!\left(
      \hat{D}^{-1/2}\hat{A}\hat{D}^{-1/2}\bd{H}^{(k)} \bd{W}^{(k)}
    \right)
  \]

  \begin{itemize}
    \item Degree correction appears on both source and target nodes
    \item Symmetric normalization is the classic GCN choice\slideref{kipf2017}
  \end{itemize}
\end{frame}

% TODO: Add a block diagram showing separate self-representation and neighborhood summary before concatenation.
\begin{frame}{Message-Passing Variants: GraphSAGE\slideref{hamilton2017}}
  \small
  \[
    d_{v} = \sum_{u=1}^{N} A_{vu}, \qquad
    D = \operatorname{diag}(d_{1}, \dots, d_{N})
  \]

  \begin{itemize}
    \item Explicit self and neighbor concatenation
  \end{itemize}
  \[
    \bd{m}_{v}^{(k)} =
    \sum_{u \in \mathcal N(v)}
    \frac{A_{vu}}{d_{v}} \,
    \bd{h}_{u}^{(k)}
  \]

  \textcolor{ntuablue}{Or in matrix notation}
  \[
    \bd{M}^{(k)} = D^{-1}A\bd{H}^{(k)}
  \]

  \textcolor{ntuablue}{Node-wise update}
  \[
    \bd{h}_{v}^{(k+1)} =
    \sigma\!\left(
      \left[
        \bd{h}_{v}^{(k)} \,;\,
        \bd{m}_{v}^{(k)}
      \right]
      \bd{W}^{(k)}
    \right)
  \]

  \textcolor{ntuablue}{Or in matrix notation}
  \[
    \bd{H}^{(k+1)} =
    \sigma\!\left(
      \left[
        \bd{H}^{(k)} \;\;
        \bd{M}^{(k)}
      \right]
      \bd{W}^{(k)}
    \right)
  \]

  \begin{itemize}
    \item \(\bd{m}_{v}^{(k)}\): average of neighbor embeddings
    \item Concatenate self state and neighbor summary before the learned update
\end{itemize}
\end{frame}

% TODO: Add an attention-weighted neighborhood plot or a heatmap of edge attention coefficients.
\begin{frame}{Message-Passing Variants: GAT Weights\slideref{velickovic2018}}
  \small
  \begin{itemize}
    \item Learned attention weights
  \end{itemize}
  \[
    \bd{Z}^{(k)} = \bd{H}^{(k)} \bd{W}^{(k)}, \qquad
    \bd{z}_{v}^{(k)} = \bd{h}_{v}^{(k)} \bd{W}^{(k)}
  \]

  \vspace{-0.18cm}
  \[
    e_{vu}^{(k)} =
    \operatorname{LeakyReLU}\!\left(
      \bd{a}^{\top}
      \begin{bmatrix}
        \bd{z}_{v}^{(k)} \\
        \bd{z}_{u}^{(k)}
      \end{bmatrix}
    \right)
  \]

  \vspace{-0.18cm}
  \[
    \alpha_{vu}^{(k)} =
    \frac{\exp\!\left(e_{vu}^{(k)}\right)}
    {\sum_{j \in \hat{\mathcal N}(v)} \exp\!\left(e_{vj}^{(k)}\right)}
  \]

  \begin{itemize}
    \item \(e_{vu}^{(k)}\): raw score for neighbor \(u\)
    \item \(\alpha_{vu}^{(k)}\): normalized attention weight around node \(v\)
  \end{itemize}
\end{frame}

\begin{frame}{Message-Passing Variants: GAT Update\slideref{velickovic2018}}
  \small
  \begin{itemize}
    \item Learned attention weights in the update
  \end{itemize}

  \textcolor{ntuablue}{Node-wise update}
  \[
    \bd{h}_{v}^{(k+1)} =
    \sigma\!\left(
      \sum_{u \in \hat{\mathcal N}(v)}
      \alpha_{vu}^{(k)} \bd{z}_{u}^{(k)}
    \right)
  \]

  \textcolor{ntuablue}{Or in matrix notation}
  \[
    \bd{H}^{(k+1)} = \sigma\!\left(\Att^{(k)} \bd{Z}^{(k)}\right),
    \qquad
    {\left(\Att^{(k)}\right)}_{vu} = \alpha_{vu}^{(k)}
  \]

  \begin{itemize}
    \item \(\Att^{(k)}\): matrix of attention coefficients
    \item Unlike GCN, the weights are learned from the node features
  \end{itemize}
\end{frame}

\sectiondivider{Prediction Levels}

% TODO: Add a visual that contrasts node outputs, edge scores, and graph readout in the same graph.
\begin{frame}{Prediction Levels: Node and Edge}
  \small
  \textcolor{ntuablue}{Node-level}
  \begin{itemize}
    \item Use the final embedding of one node
  \end{itemize}
  \[
    \hat{y}_{v} = g\!\left(\bd{h}_{v}^{(L)}\right)
  \]

  \textcolor{ntuablue}{Edge-level}
  \begin{itemize}
    \item Combine the two endpoint embeddings
  \end{itemize}
  \[
    \hat{y}_{uv} = g\!\left(\bd{h}_{u}^{(L)}, \bd{h}_{v}^{(L)}\right)
  \]
\end{frame}

\begin{frame}{Prediction Levels: Graph-Level}
  \small
  \begin{itemize}
    \item Pool node embeddings into one graph representation
  \end{itemize}
  \[
    \bd{h}_{G} =
    \operatorname{READOUT}\!\left(\{\bd{h}_{v}^{(L)} : v \in V\}\right)
  \]

  \begin{itemize}
    \item Common choices
  \end{itemize}
  \[
    \bd{h}_{G} =
    \sum_{v \in V} \bd{h}_{v}^{(L)}
    \qquad \text{or} \qquad
    \bd{h}_{G} =
    \frac{1}{N}\sum_{v \in V} \bd{h}_{v}^{(L)}
  \]

  \begin{itemize}
    \item Then apply a graph-level decoder: \(\hat{y}_{G} = g(\bd{h}_{G})\)
  \end{itemize}
\end{frame}

\begin{frame}{Summary}
  \[
    \bd{H}^{(k+1)} =
    \sigma\!\left(P^{(k)} \bd{H}^{(k)} \bd{W}^{(k)}\right)
  \]

  \begin{itemize}
    \item Weighted averaging: \(P = \hat{D}^{-1}\hat{A}\)
    \item GCN: \(P = \hat{D}^{-1/2}\hat{A}\hat{D}^{-1/2}\)
    \item GraphSAGE: combine self state with an explicit neighbor summary
    \item GAT: replace fixed weights with learned attention coefficients
    \item Same message-passing template, different weighting rule
  \end{itemize}
\end{frame}

\begin{frame}{Reference}
  \small
  Sanchez-Lengeling, Reif, Pearce, Wiltschko.

  \vspace{0.15cm}
  \emph{A Gentle Introduction to Graph Neural Networks}. Distill, 2021.

  \vspace{0.2cm}
  \href{https://distill.pub/2021/understanding-gnns/}{https://distill.pub/2021/understanding-gnns/}
\end{frame}

\end{document}
